% \chapter{Phân công nhiệm vụ}
% \section{Nhận xét quá trình làm việc}
% Trong quá trình làm việc, các thành viên trong nhóm đã tích cực nghiên cứu, hoàn thành công việc theo đúng tiến độ đề ra. Nhóm cũng thực hiện báo cáo tiến độ cho giảng viên hướng dẫn theo lịch trình một cách đầy đủ và chi tiết để nhận được những nhận xét, góp ý kịp thời.

% Về mặt phân công, trong suốt quá trình làm việc, các thành viên đều có sự tham gia ít nhiều đối với tất cả các phần của Đồ án. Điều này giúp các thành viên hiểu và nắm rõ toàn bộ nội dung của đồ án, cũng như có thể trực tiếp tham gia hỗ trợ lẫn nhau và dễ dàng hơn trong việc phối hợp. Bên cạnh đó, để có để đạt được hiệu quả tốt nhất thì việc phân công vẫn dựa trên điểm mạnh và điểm yếu của từng thành thành viên, nhưng vẫn đảm bảo tính hợp lý về mặt thời gian và công bằng về khối lượng công việc.
% \section{Chi tiết phân công nhiệm vụ}
% Chi tiết phân công nhiên vụ cho từng thành viên được thể hiện trong bảng \ref{tab:phan_cong}.

% \begin{table}[!htp]
%     \centering
%     \caption{Bảng chi tiết phân công nhiệm vụ}
%     \begin{tabular}{|c|c|}
%     \hline
%         \textbf{Thành viên} & \textbf{Phân công nhiệm vụ}\\
%         \hline
%         \textbf{Hà Trung Kiên} & \makecell[l]{
%         Nghiên cứu các công trình liên quan.\\
%         Nghiên cứu các kiến thức nền tảng.\\
%         Phân tích và thiết kế các yêu cầu hệ thống.\\
%         Hiện thực các microservice về User, Project,...\\
%         Nghiên cứu các phương án cần thiết. \\
%         Triển khai hệ thống.
%         \\
%         }\\
%         \hline
%         \textbf{Lương Hồng Tiến Đạt} &  \makecell[l]{Nghiên cứu các kiến thức nền tảng.\\Nghiên cứu các công trình liên quan.\\Phân tích và thiết kế các yêu cầu hệ thống.\\ Thiết kế Database và kiến trúc hệ thống. \\ Hiện thực các microservice về Suggestion, Forum,... \\
%         Triển khai hệ thống.
%         \\}\\
%         \hline
%         \textbf{Hà Phan Thiên Phú} &  \makecell[l]{Nghiên cứu các kiến thức nền tảng.\\Nghiên cứu các công trình liên quan.\\
%          Phân tích và thiết kế các yêu cầu hệ thống.\\
%          Thiết kế giao diện cho hệ thống.\\
%         Hiện thực phần frontend cho hệ thống.\\
%         Triển khai hệ thống.
%          }\\
%         \hline
%     \end{tabular}
%     \label{tab:phan_cong}
% \end{table}